\chapter{Background and Mathematical Model}
\label{chap:background}


We will use a unified notation for concepts across \todo{ref all papers}.



% ---------------- Template Examples ---------------------
% \section{Notation and array accesses}
% Use inline mathematics for short expressions, such as a loop index
% $\vect{i}\in\Z^n$, and displayed equations for derivations. An affine access
% function can be written as
% \begin{equation}
%   \vect{a}(\vect{i}) = H\vect{i} + \vect{c},
%   \qquad H\in\Z^{m\times n},\quad\vect{c}\in\Z^m.
%   \label{eq:affine-access}
% \end{equation}
% Here, $\vect{a}$ denotes array subscripts rather than a byte address; converting
% subscripts to an address also depends on element size and storage layout.
%
% \begin{definition}[Temporal and spatial locality]
% Temporal locality concerns accesses to data used recently. Spatial locality
% concerns accesses to nearby memory locations, which may share a cache line.
% \end{definition}
%
% \begin{table}[htbp]
%   \centering
%   \caption{Notation used in the illustrative model.}
%   \label{tab:notation}
%   \begin{tabularx}{\textwidth}{lY}
%     \toprule
%     Symbol & Meaning \\
%     \midrule
%     $\vect{i}$ & Integer iteration vector \\
%     $H$, $\vect{c}$ & Array-access matrix and constant offset \\
%     $L$ & Cache-line size in bytes \\
%     $C$ & Available cache capacity in bytes \\
%     $e$ & Element size in bytes \\
%     $N$, $B$ & Matrix dimension and tile dimension, respectively \\
%     \bottomrule
%   \end{tabularx}
% \end{table}
%
% \section{An illustrative cost model}
% For demonstration, average access time may be expressed as
% \begin{align}
%   T_{\mathrm{avg}}
%     &= T_{\mathrm{hit}} + p_{\mathrm{miss}}T_{\mathrm{penalty}},
%     \label{eq:access-time}\\
%   p_{\mathrm{miss}}
%     &= \frac{M}{A}, \qquad A>0,
%     \label{eq:miss-rate}
% \end{align}
% where $M$ is the number of misses and $A$ is the number of accesses for the
% cache level being modelled. The miss penalty is additional to the hit time.
% \Cref{eq:access-time} is a simplified explanatory model; it does not
% capture overlapping misses, prefetching, or a complete cache hierarchy.
%
% For a rough capacity estimate, three dense $B\times B$ tiles of elements of size
% $e$ require
% \begin{equation}
%   W(B) = 3eB^2 \leq \alpha C, \qquad 0<\alpha<1.
%   \label{eq:tile-capacity}
% \end{equation}
% The factor $\alpha$ leaves room for other data. This estimate alone does not
% account for cache conflicts or establish an optimal tile size.
%
% \section{Transformation notation}
% A matrix representation is convenient when discussing loop transformations.
% For example, interchanging two loop dimensions can be represented by
% \begin{equation}
%   \vect{i}' = U\vect{i}, \qquad
%   U = \begin{bmatrix}0 & 1 \\ 1 & 0\end{bmatrix}.
%   \label{eq:interchange}
% \end{equation}
% \todo{Add dependence constraints and explain the legality conditions used by
% your selected papers; see \citet{wolfLam1991parallel} for the transformation-theory
% paper in the reading list.}
